\documentclass[11pt,a4paper]{article}

\usepackage[a4paper,margin=1in]{geometry}
\usepackage{amsmath}
\usepackage{array}
\usepackage{booktabs}
\usepackage{enumitem}
\usepackage{tabularx}
\usepackage{tikz}
\usepackage{fancyhdr}
\usepackage[T1]{fontenc}
\usepackage{lmodern}

\usetikzlibrary{arrows.meta,positioning}

\newcommand{\universityname}{UNIVERSITY OF DERBY}
\newcommand{\college}{College of Science and Engineering}
\newcommand{\examtitle}{Practice Paper B}
\newcommand{\moduletitle}{OPERATIONAL RESEARCH}
\newcommand{\modulecode}{5CM522}
\newcommand{\examdate}{REVISION MATERIAL}
\newcommand{\timeallowed}{4 HOURS}

\newcommand{\continuednotice}{%
    \vfill
    \begin{center}
    \textit{Continued on Next Page}
    \end{center}
    \newpage
}

\newcommand{\markalloc}[1]{\nobreak\hfill\mbox{(#1 Marks)}}
\newcommand{\questiontotal}[1]{\begin{flushright}\textbf{Total #1 Marks}\end{flushright}}

\pagestyle{fancy}
\fancyhf{}
\fancyfoot[C]{\thepage}
\renewcommand{\headrulewidth}{0pt}

\setlength{\parindent}{0pt}
\setlength{\parskip}{0.4em}
\setlist[enumerate,1]{label=(\alph*), leftmargin=1.1cm, itemsep=0.55em, topsep=0.4em}

\begin{document}

\begin{center}
{\large \textbf{\universityname}}\\
{\college}
\end{center}

\vspace{0.5em}

\begin{center}
{\Large \textbf{\examtitle}}\\[0.5em]
\rule{\textwidth}{0.4pt}\\[0.6em]
{\large \textbf{\moduletitle}}\\
{\large \textbf{\modulecode}}\\[0.4em]
\textbf{DATE:} \examdate\\
\textbf{TIME ALLOWED:} \timeallowed\\[0.3em]
\rule{\textwidth}{0.4pt}
\end{center}

\textbf{Instructions to Candidates}

You are required to answer a total of 4 questions.

However, you may attempt all 6 and your best 4 will be considered.

This is a practice paper designed to support preparation for the timed in-class test.

You may use Excel or Python, but your workings, files and interpretations must be clear and well labelled.

You may access code, notes and module files while practising.

You must not use AI tools if you want the practice to reflect the conditions of the real in-class test.

\newpage

\textbf{1.}

Silverforge Engineering is studying two scheduling problems.

\textbf{Part A: Flow Shop Scheduling}

All five jobs must be processed first on Machine 1 and then on Machine 2.

\begin{center}
\begin{tabular}{lcc}
\toprule
Job & Machine 1 Time & Machine 2 Time \\
\midrule
$J_1$ & 7 & 5 \\
$J_2$ & 3 & 8 \\
$J_3$ & 6 & 4 \\
$J_4$ & 2 & 9 \\
$J_5$ & 5 & 6 \\
\bottomrule
\end{tabular}
\end{center}

\textbf{Part B: Job Shop Scheduling}

Each job has its own machine order and must follow that order exactly.

\begin{center}
\begin{tabular}{lcc}
\toprule
Job & Sequence on Machines & Processing Times \\
\midrule
$A$ & $M_1, M_3, M_2$ & $4, 5, 3$ \\
$B$ & $M_2, M_1, M_3$ & $3, 4, 6$ \\
$C$ & $M_3, M_2, M_1$ & $5, 2, 4$ \\
$D$ & $M_1, M_2, M_3$ & $2, 6, 3$ \\
\bottomrule
\end{tabular}
\end{center}

\begin{enumerate}
\item Use Johnson's Rule to determine an optimal production sequence for the flow shop problem in Part A. Show the main steps. \markalloc{6}

\item Draw a Gantt chart for your sequence and calculate the makespan. \markalloc{5}

\item Compare your Johnson schedule with the original order $J_1,J_2,J_3,J_4,J_5$ and comment on the improvement. \markalloc{4}

\item Solve the job shop scheduling problem in Part B and state an optimal schedule and the corresponding makespan. You may present the result as a table of operation start and finish times or as a Gantt chart. \markalloc{6}

\item Suppose job $C$ must finish by time 15. Explain how the model changes and discuss the possible effect on feasibility and makespan. \markalloc{4}
\end{enumerate}

\questiontotal{25}

\continuednotice

\textbf{2.}

County Response wants to choose a limited number of patrol bases.
The aim is not to cover every district, but to maximise the total population covered while opening at most three bases.

\begin{center}
\begin{tabular}{lccc}
\toprule
Base & Cost & Population Covered & Districts Covered \\
\midrule
$A$ & 14 & 18{,}000 & $\{1,2,4\}$ \\
$B$ & 11 & 15{,}000 & $\{1,3\}$ \\
$C$ & 13 & 21{,}000 & $\{2,3,5\}$ \\
$D$ & 10 & 12{,}000 & $\{4,5\}$ \\
$E$ & 12 & 16{,}000 & $\{2,5,6\}$ \\
$F$ & 9  & 14{,}000 & $\{3,4,6\}$ \\
\bottomrule
\end{tabular}
\end{center}

The population living in each district is given below.

\begin{center}
\begin{tabular}{lcccccc}
\toprule
District & 1 & 2 & 3 & 4 & 5 & 6 \\
\midrule
Population & 9{,}000 & 12{,}000 & 14{,}000 & 10{,}000 & 11{,}000 & 8{,}000 \\
\bottomrule
\end{tabular}
\end{center}

The budget is \pounds31 and at most 3 bases may be opened.

\begin{enumerate}
\item Formulate a binary integer programming model to maximise covered population subject to the budget and base limit. Clearly define any coverage variables you use. \markalloc{9}

\item Solve the model and identify which bases should be opened and the resulting total covered population. \markalloc{6}

\item Explain whether some districts may remain uncovered in an optimal solution and why this is acceptable in this model but not in a set covering model. \markalloc{4}

\item Modify the model so that district 6 must be covered and bases $B$ and $F$ cannot both be chosen. State the revised model changes clearly. \markalloc{6}
\end{enumerate}

\questiontotal{25}

\newpage

\textbf{3.}

A university library help desk is modelled as a queueing system.
Students arrive randomly and service times are assumed to follow an exponential distribution.
You should adapt the SimPy structures used in class, but tailor the model to this setting.

Use the following assumptions:

\begin{center}
\begin{tabular}{lc}
\toprule
Quantity & Value \\
\midrule
Simulation length & 360 minutes \\
Average service time & 7 minutes \\
Baseline arrival rate & 10 students per hour \\
Peak arrival rate & 15 students per hour \\
Service target & 97\% wait less than 8 minutes \\
Priority users & 10\% of arrivals \\
Suggested random seed & 24 \\
\bottomrule
\end{tabular}
\end{center}

\begin{enumerate}
\item Describe the main elements of a discrete-event simulation model for this system: entities, resources, events, stochastic inputs and performance measures. \markalloc{5}

\item Adapt a SimPy model to represent the baseline system with 2 advisers and an arrival rate of 10 students per hour. Report the total served, average wait, maximum wait and percentage waiting less than 8 minutes. \markalloc{6}

\item Investigate the peak-period scenario with 2, 3 and 4 advisers. Present your results in a clearly labelled table and determine the minimum number of advisers needed to meet the service target. \markalloc{4}

\item Extend the model so that 10\% of users are priority users who should be served ahead of standard users. Explain how the model should be changed and discuss the likely effect on standard-user waits. \markalloc{6}

\item Explain why a single run is not enough to support a staffing decision and describe how you would use replications and/or a warm-up period to improve the analysis. \markalloc{4}
\end{enumerate}

\questiontotal{25}

\continuednotice

\textbf{4.}

The operations analyst at Broadleaf Homeware has collected the following quarterly demand data.

\begin{center}
\begin{tabular}{cccc}
\toprule
Period & Year & Quarter & Demand \\
\midrule
1  & 2019 & Q1 & 420 \\
2  & 2019 & Q2 & 500 \\
3  & 2019 & Q3 & 470 \\
4  & 2019 & Q4 & 560 \\
5  & 2020 & Q1 & 450 \\
6  & 2020 & Q2 & 536 \\
7  & 2020 & Q3 & 500 \\
8  & 2020 & Q4 & 600 \\
9  & 2021 & Q1 & 485 \\
10 & 2021 & Q2 & 575 \\
11 & 2021 & Q3 & 535 \\
12 & 2021 & Q4 & 645 \\
13 & 2022 & Q1 & 522 \\
14 & 2022 & Q2 & 618 \\
15 & 2022 & Q3 & 572 \\
16 & 2022 & Q4 & 694 \\
\bottomrule
\end{tabular}
\end{center}

\begin{enumerate}
\item Using periods 1--12 only, calculate suitable centred 4-quarter moving averages and estimate the trend values where possible. Plot the original data together with the estimated trend. \markalloc{8}

\item Using periods 1--12 only, calculate seasonal indices using both additive and multiplicative decomposition models. \markalloc{7}

\item Use the 2022 data (periods 13--16) as a validation set. Generate forecasts for these four periods using both models and compare their accuracy using one suitable measure such as MAD, MSE or MAPE. \markalloc{6}

\item Using the better model, produce forecasts for the four quarters of 2023 and justify your final model choice. \markalloc{4}
\end{enumerate}

\questiontotal{25}

\newpage

\textbf{5.}

Harbour Plastics makes two products each week: casings and inserts.

Each casing requires 4 moulding hours and 1 assembly hour.
Each insert requires 2 moulding hours and 3 assembly hours.
At most 100 moulding hours and 90 assembly hours are available.
No more than 24 inserts can be sold each week.

The weekly contribution is \pounds40 per casing and \pounds34 per insert.

\begin{enumerate}
\item Formulate a linear programming model for this problem. \markalloc{6}

\item Solve the model to determine the production plan that maximises weekly contribution. \markalloc{6}

\item State the binding constraints and the unused capacity in each department. \markalloc{4}

\item A distributor now requires at least 25\% of total output to be inserts. Reformulate the model and determine whether the original solution remains feasible. \markalloc{5}

\item If up to 12 extra assembly hours can be purchased at \pounds3 per hour for one week only, determine whether they should be bought. \markalloc{4}
\end{enumerate}

\questiontotal{25}

\continuednotice

\textbf{6.}

Northline Distribution must ship goods from three warehouses to four stores at minimum cost.

\begin{center}
\begin{tabular}{lcccc}
\toprule
Warehouse & $S_1$ & $S_2$ & $S_3$ & $S_4$ \\
\midrule
$W_1$ & 6 & 8 & 10 & 7 \\
$W_2$ & 9 & 5 & 7 & 6 \\
$W_3$ & 4 & 7 & 6 & 8 \\
\bottomrule
\end{tabular}
\end{center}

\begin{center}
\begin{tabular}{lc}
\toprule
Warehouse & Supply \\
\midrule
$W_1$ & 80 \\
$W_2$ & 100 \\
$W_3$ & 90 \\
\bottomrule
\end{tabular}
\hspace{1.5cm}
\begin{tabular}{lc}
\toprule
Store & Demand \\
\midrule
$S_1$ & 70 \\
$S_2$ & 60 \\
$S_3$ & 80 \\
$S_4$ & 60 \\
\bottomrule
\end{tabular}
\end{center}

\begin{enumerate}
\item Formulate a transportation model for this problem, clearly defining the decision variables, objective function and constraints. \markalloc{7}

\item Solve the model and determine the optimal shipping plan and total transportation cost. Report only the non-zero shipments. \markalloc{7}

\item Resolve the model if warehouse $W_2$ cannot supply store $S_2$, and at least 20 units must be shipped from $W_1$ to $S_4$. Comment on the impact on total cost. \markalloc{6}

\item Explain one realistic extension that would make this transportation model more useful in practice. \markalloc{5}
\end{enumerate}

\questiontotal{25}

\vfill

\begin{center}
\textbf{End of Paper}
\end{center}

\end{document}
